\documentclass[../main.tex]{subfiles}

\begin{document}
\subsection{Consecuencia de la fórmula integral de Cauchy y Principio del Módulo Máximo}
\classdate{2026-04-23}

Veremos una consecuencia de la fórmula integral de Cauchy (FIC).

\begin{remark}
    Recordemos la condición

    \begin{equation*}
        \lim_{z \to z_{0}} f(z)\bigl(z - z_{0}\bigr) = 0
    \end{equation*}

    que pedimos en el \zcref{lemma:lem-goursat-v2},
    para el cual reservamos un nombre especial: singularidad removible.
\end{remark}

\begin{definition}
    Sean \(\omega \subseteq \mathbb{C}\) abierto, \(z_{0}\in  \Omega\) y \(f \colon
    \Omega\setminus \{z_{0}\} \to \mathbb{C}\) analítica. Decimos que \(z_{0}\) es una
    singularidad removible de \(f\) si y solo si

    \begin{equation*}
        \lim_{z \to z_{0}} f(z)\bigl(z - z_{0}\bigr) = 0.
    \end{equation*}
\end{definition}

\begin{proposition}
    Sean \(\Omega \subseteq \mathbb{C}\) abierto, \(z_{0}\in \Omega\) y \(f \colon
    \Omega\setminus \{z_{0}\} \to \mathbb{C}\) una función analítica. \(z_{0}\) es una
    singularidad removible de \(f\) si y solo si \(\exists \tilde{f} \colon \Omega \to
    \mathbb{C}\) tal que \(\tilde{f}(z) = f(z)\; \forall z\in \Omega\setminus
    \{z_{0}\}\).
\end{proposition}

\begin{proof}
    \begin{itemize}
        \item Si \(\exists \tilde{f} \implies z_{0}\) es una singularidad removible

            Supongamos que \(\exists \tilde{f}\colon \Omega \to
            \mathbb{C}\) analítica tal que
            \(\tilde{f}(z) = f(z)\; \forall z\in \Omega\setminus \{z_{0}\}\). Entonces,
            \(\tilde{f}\) es continua en \(z_{0}\) y por lo tanto,

            \begin{equation*}
                \lim_{z \to z_{0}} f(z)(z - z_{0}) = \lim_{z \to z_{0}}
                \tilde{f}(z)(z - z_{0}) =
                \tilde{f}(z_{0})(0) = 0.
            \end{equation*}

            por lo que \(z_{0}\) es una singularidad removible.

        \item Si \(z_{0}\) es una singularidad removible, entonces \(\exists \tilde{f}\)

            Supongamos que \(z_{0}\) es una singularidad removible, es decir,

            \begin{equation*}
                \lim_{z \to z_{0}} f(z)(z - z_{0}) = 0.
            \end{equation*}

            Usando la FIC definimos

            \begin{equation*}
                \tilde{f}(z) = \dfrac{1}{2\pi
                i}\int_{\gamma}^{}\dfrac{f(w)}{w - z}\odif{z},
            \end{equation*}

            donde \(\gamma  \colon [a, b] \to \Omega\) es una curva cerrada
            simple suave por pedazos
            (que rodea una vez a \(z_{0}\)). De esta manera, \(\forall z\in
            \Omega\), \(z \neq
            z_{0}, z\in \gamma, \tilde{f}(z) = f(z)\) y por el
            \zcref{lemma:lem-aux}, \(\tilde{f}\) es
            analítica y está definida en \(z_{0}\).
    \end{itemize}
\end{proof}

\begin{theorem}[Principio del módulo máximo]\label{thm:maximumModulusPrinciple}
    Sea \(\Omega\) una región y \(f \colon \Omega \to \mathbb{C}\) una función
    analítica. Si \(\abs{f(z)}\) alcanza su valor máximo en \(z_{0}\in \Omega\),
    entonces \(f\) es constante.
\end{theorem}

\begin{proof}
    Primero mostraremos que \(f\) es constante en un disco abierto centrado en
    \(z_{0}\). Si \(\gamma \colon [0, 2\pi] \to D\), \(\gamma(t) = z_{0} + r
    \mathrm{e}^{it}\), con \( r < \mathrm{radio}(D)\). \(\gamma\) es una
    curva cerrada simple
    suave por pedazos, que da una vuelta a \(z_{0}\). Por la FIC,

    \begin{align*}
        f(z_{0}) &= \dfrac{1}{2\pi i}\int_{\gamma}^{}\dfrac{f(z)}{z - z_{0}}\odif{z},\\
        &= \dfrac{1}{2\pi i}\int_{0}^{2\pi}\dfrac{f(\gamma(t))}{\gamma(t) -
        z_{0}}\gamma^{\prime}(t)\odif{t},\\
        &= \dfrac{1}{2\pi i}\int_{0}^{2\pi}\dfrac{f(z_{0} + r \mathrm{e}^{it})}{r
        \mathrm{e}^{it}}ri \mathrm{e}^{it}\odif{t},\\
        &= \dfrac{1}{2\pi}\int_{0}^{2\pi}f(z_{0} + r \mathrm{e}^{it})\odif{t}.
    \end{align*}

    Tomemos módulos

    \begin{equation*}
        \abs{f(z_{0})} \leq \dfrac{1}{2\pi}\int_{0}^{2\pi}\abs{f(z_{0} + r
        \mathrm{e}^{it})}\odif{t} \leq
        \dfrac{1}{2\pi}\int_{0}^{2\pi}\abs{f(z_{0})}\odif{t} =
        \abs{f(z_{0})},
    \end{equation*}

    pero entonces

    \begin{equation*}
        \abs{f(z_{0})} = \dfrac{1}{2\pi} \int_{0}^{2\pi}\abs{f(z_{0} + r
        \mathrm{e}^{it})}\odif{t} = 0.
    \end{equation*}

    Como el integrando es continuo y mayor o igual a 0, tenemos que \(\abs{f(z_{0})} =
    \abs{f(z_{0} + r \mathrm{e}^{it})}\). Por lo que \(\abs{f}\) es constante en la
    circunferencia de radio \(r\) centrada en \(z_{0}\), donde \(r <
    \mathrm{radio}(D)\).

    Denotamos por \(\gamma_{r}\) a esta circunferencia. Como \(D\setminus \{z_{0}\} =
    \bigcup_{r < \mathrm{radio}(D)}\gamma_{r}\), concluimos que \(\abs{f}\)
    es constante en
    \(D\setminus \{z_{0}\}\) (la constante \(\abs{f(z_{0})}\)). Por lo que
    \(\abs{f}\) es
    constante en \(D\).

    Veamos que esto implica que \(f\) es constante en \(D\). Escribimos \(f = u + iv\),
    \(\abs{f}^{2} = u^{2} + v^{2} = C = \abs{f(z_{0})}^{2}\). Si
    \(\abs{f(z_{0})}^{2} \geq
    0\), entonces \(\abs{f(z)} \leq \abs{f(z_{0})} = 0\), lo cual implica que
    \(f(z) = 0\;
    \forall z\in \mathbb{C}\), por lo que \(f\) es constante.

    Supongamos que \(C > 0\). Derivamos parcialmente \(u^{2} + v^{2} = C\)
    respecto a \(x\),

    \begin{equation*}
        \pdv*{\bigl(u^{2} + v^{2}\bigr)}{x} = \pdv{c}{k} = 0.
    \end{equation*}

    Es decir,

    \begin{align*}
        2u \pdv{u}{x} + 2v \pdv{v}{x} &= 0,\\
        u \pdv{u}{x} + v \pdv{v}{x} &= 0.
    \end{align*}

    Análogamente, derivamos respecto a \(y\),

    \begin{equation*}
        u \pdv{u}{y} + v \pdv{v}{y} = 0.
    \end{equation*}

    Tenemos el sistema

    \begin{equation*}
        \begin{dcases}
            u \pdv{u}{x} + v \pdv{v}{x} = 0,\\
            u \pdv{u}{y} + v \pdv{v}{y} = 0.
        \end{dcases}
    \end{equation*}

    Como \(f\) es analítica, satisface Cauchy-Riemann. Entonces podemos reescribir el
    sistema, usando \(\pdv{u}!{x} = \pdv{v}!{y}\) y \(\pdv{u}!{y} = - \pdv{v}!{x}\),

    \begin{equation*}
        \begin{dcases}
            u\pdv{u}{x} - v\pdv{u}{y} = 0,\\
            v\pdv{u}{x} + u\pdv{u}{y} = 0.
        \end{dcases}
    \end{equation*}

    Esto lo podríamos escribir de forma matricial,

    \begin{equation*}
        \begin{pNiceMatrix}
            u & -v\\
            v & u\\
        \end{pNiceMatrix}
        \begin{pNiceMatrix}
            \pdv{u}{x}\\
            \pdv{v}{y}\\
        \end{pNiceMatrix} =
        \begin{pNiceMatrix}
            0\\
            0\\
        \end{pNiceMatrix}.
    \end{equation*}

    Tal que,

    \begin{equation*}
        \det
        \begin{pNiceMatrix}
            u & -v\\
            v & u\\
        \end{pNiceMatrix}
        = u^{2} + v^{2} = c > 0.
    \end{equation*}

    Por lo que
    \begin{equation*}
        \pdv{u}{x} = 0\quad \wedge \quad \pdv{v}{y} = 0.
    \end{equation*}

    Pero por Cauchy-Riemann, \(\pdv{v}{x} = 0 = \pdv{v}{y}\).

    Como \(D\) es conexo, tenemos que \(u\) y \(v\) son constantes. Por lo que \(f\) es
    constante en \(D\).

    Sea \(A = \bigl\lbrace z\in \Omega \mid f(z) = f(z_{0}) \bigr\rbrace\).
    Por lo anterior,
    \(A\) es abierto. Además, \(A \neq \emptyset\) pues \(z_{0}\in  A\) y
    \(A\) es cerrado
    pues \(A = f^{-1}\bigl(\{f(z_{0})\}\bigr)\). Por lo que \(A\) es abierto,
    cerrado y no
    vacío en \(\Omega\), el cual es conexo. Y además, \(A = \Omega\). Por lo
    tanto, \(f\) es
    constante en \(\Omega\).
\end{proof}

\classdate{2026-04-24}

\begin{corollary}
    Sean \(\Omega\) una región acotada de \(\mathbb{C}\) y \(f \colon \overline{\Omega}
    \to \mathbb{C}\) continua en \(\overline{\Omega}\).
\end{corollary}

\begin{proof}
    Como \(\Omega\) es acotado, entonces \(\overline{\Omega}\) es acotado. Por lo que
    \(\overline{\Omega}\) es cerrado y acotado, i.e. \(\overline{\Omega}\) es compacto.
    Notemos que \(\abs{f} \colon \overline{\Omega} \to  \mathbb{R}\) es una función
    definida en un compacto, tal que \(\abs{f}\) alcanza su máximo en
    \(\overline{\Omega}\).

    Ahora si la función alcanza su máximo en \(\partial \Omega\), ya no hay nada que
    demostrar. Pero si el máximo se alcanza en \(\Omega\), por el
    \zcref{thm:maximumModulusPrinciple},
    \(f\) es constante en \(\Omega\). Afirmamos que \(f\) es constante en
    \(\partial\Omega\)
    y que ambas constantes coinciden. Y supongamos que \(f \equiv c\) en \(\Omega\). Sea
    \(z\in \partial\Omega\), entonces existe \(\bigl\lbrace z_{n} \bigr\rbrace_{n\in
    \mathbb{N}} \subseteq \Omega\) tal que \(z_{n} \to z\). Como \(f\) es continua en
    \(\overline{\Omega}\), \(f(z_{n}) \to f(z)\), \(c \to f(z)\) y \(f \equiv c\) en
    \(\overline{\Omega}\). Por lo que \(\exists z_{0}\in \partial\Omega\) tal que
    \(\underbrace{\abs{f(z)}}_{= c} \leq \underbrace{\abs{f(z_{0})}}_{= c}\;
        \forall z\in
    \overline{\Omega}\).
\end{proof}

\begin{lemma}[de Schwarz]
    Sean \(D = B_{1}(0)\) el disco unitario y \(f \colon D \to \mathbb{C}\) una función
    analítica tal que \(\abs{f(z)} \leq 1\; \forall z\in D\) y \(f(0) = 0\). Entonces,

    \begin{enumerate}
        \item \(\abs{f(z)} \leq \abs{z}\; \forall z\in D\).
        \item \(\abs{f^{\prime}(0)} \leq 1\).
    \end{enumerate}

    Si además se cumple que \(\abs{f(z_{0})} = \abs{z_{0}}\) para algún \(z_{0}\in
    D\setminus \{0\}\) o bien \(\abs{f^{\prime}(0)} = 1\), entonces \(f(z) =
    cz\) para algún
    \(z\in \mathbb{C}\) tal que \(\abs{c} = 1\).
\end{lemma}

\begin{proof}
    Definimos

    \begin{equation*}
        g(z) =
        \begin{dcases}
            \dfrac{f(z)}{z}, & z \neq 0\\
            f^{\prime}(0), & z = 0.
        \end{dcases}
    \end{equation*}

    Como \(f\) es analítica en \(D\setminus \{0\}\). Además para \(z \neq
    0\), \(g(z) z =
    f(z)\). Por lo que,

    \begin{equation*}
        \lim_{z \to 0} g(z) z = \lim_{z \to 0} f(z_{0}) = f(0) 0.
    \end{equation*}

    Es decir, \(0\) es una singularidad removible de \(g\). Y \(\exists
        \tilde{g} \colon D
    \to \mathbb{C}\) analítica tal que \(\tilde{g}(z) = g(z)\; \forall z \neq 0\).
\end{proof}

Por otro lado, como \(f(0) = 0\), para \(z \neq 0\),

\begin{align*}
    g(z) &= \dfrac{f(z)}{z} = \dfrac{f(z) - f(0)}{z - 0},\\
    \implies\;\lim_{z \to 0} g(z) &= \lim_{z \to 0} \dfrac{f(z) - f(0)}{z -
    0} = f^{\prime}(0).
\end{align*}

Por lo que \(g\) es continua en 0 y \(g = \tilde{g}\) en todo \(D\). Entonces, \(g\) es
analítica en todo el disco.

Consideremos ahora \(r < 1\) y \(B_{r}(0) \subsetneq D\). Si \( \abs{z} = r\), entonces

\begin{equation*}
    \abs{g(z)} = \dfrac{\abs{f(z)}}{\abs{z}} \leq \dfrac{1}{r}.
\end{equation*}

Por el corolario anterior, \(\abs{g(z)} \leq \tfrac{1}{r}\; \forall z\in
\overline{B_{r}(0)}\).

Tomando \( r \to 1^{-}\), tenemos que \(\abs{g(z)} \leq 1\; \forall z\in D\). En
particular,

\begin{equation*}
    \abs{f(z)} \leq \abs{z}\; \forall z\in D\setminus \{0\}
\end{equation*}

y \(\abs{f^{\prime}(0)} \leq 1\). Recordando que \(f(0) = 0\).

Si \(\abs{f(z_{0})} = \abs{z_{0}}\) para algún \(z_{0}\in  D\setminus \{0\}\), entonces
\(\abs{g(z_{0})} = 1\). Por lo que el módulo máximo de \(g\) se alcanza en un punto de
\(D\). Por el principio del módulo máximo \(g\) es constante. De manera, análoga, si
\(\abs{f^{\prime}(0)} = 1\), entonces \(\abs{g(0)} = 1\) y nuevamente \(g\) es
constante. En ambos casos, \(\exists c\in \mathbb{C}\) tal que \(\abs{g(z)} = c\;
\forall z\in D\). Pero \(1 = \abs{g(z)} = c\), lo cual implica que \(\abs{c} = 1\) y
además, que \(\forall z \neq 0\), \(c = g(z) = \tfrac{f(z)}{z} \implies f(z) = cz\). Y
si \(z = 0\), \(f(z) = 0 = c \cdot 0 = cz\). Por lo tanto,

\begin{equation*}
    f(z) = cz\; \forall z\in D.
\end{equation*}
\end{document}
